\subsection{硬件作用与调度作用的成对识别}
令 $p_{B,i,k}$ 为问题二对图 $i$、核数 $k$ 的选用方案，$p_{L,i,k}$ 为问题三重新搜索后的方案。定义三次相互不同的评估：
\begin{equation}\label{eq:q3-three-evaluations}
T_B=T^B(p_B),\qquad
T_{\rm fixed}=T^{L2}(p_B),\qquad
T_{\rm selected}=T^{L2}(p_L).
\end{equation}
$T_B$ 与 $T_{\rm fixed}$ 的方案完全相同，只改变有无共享 L2；$T_{\rm fixed}$ 与 $T_{\rm selected}$ 的硬件完全相同，只改变被选方案。对同一图与同一核数，可以逐项分解
\begin{equation}\label{eq:q3-ratio-factorization}
R_{\rm hw}=\frac{T_B}{T_{\rm fixed}},\qquad
R_{\rm search}=\frac{T_{\rm fixed}}{T_{\rm selected}},\qquad
R_{\rm total}=\frac{T_B}{T_{\rm selected}}
=R_{\rm hw}R_{\rm search}.
\end{equation}
式\eqref{eq:q3-ratio-factorization}只是\emph{逐配置}的恒等式，不能把 100 张图的 $\overline R_{\rm hw}$ 与 $\overline R_{\rm search}$ 相乘后称为 $\overline R_{\rm total}$，因为平均乘积一般不等于平均值乘积。表\ref{tab:q3-main}的每一列都先在逐图层面计算，再取算术平均。三条五点加速比曲线沿用相同的无 L2 单核参考；问题三的单核 L2 比值允许大于 $1$，不能把它强制定为赛题前两问的单核点。

这一成对设计还限定了“硬件收益”的解释范围：它是在\emph{固定已选 Q2 方案}上增加 L2 的局部反事实，不是“所有可能调度下增设缓存的纯因果平均”。Q2 方案原本按无 L2 规则选优，未必是 L2 的最佳方案；问题三再搜索的收益则取决于新增候选池。只有三次评价都使用相同图、核数、原始 DDR 配置及官方代码版本，式\eqref{eq:q3-ratio-factorization}才有明确含义。

\subsection{从 FIFO 账本构造有界候选}
给定种子方案 $p_B$，先在题目默认 $C=1$ MiB、$b_L=250$ bytes/cycle 下评价，得到完整事件时间线。遍历正字节 \texttt{COPY\_IN} 事件，记录其逻辑张量、发射时 FIFO 队列、命中标志、服务池、完成时刻及相关消费子图。对重复未命中的张量，区分“首次冷读”“容量淘汰后再读”和“并发冷读”三种情况；仅有命中率数字无法指明可操作的修复位置。再把影响最终输出的较晚读入或关键链等待映射回少数相关子图，生成迁移核心或合法顺序插入候选。这样选择的仍是 $(\pi,m,\sigma)$ 的局部变化，算法没有直接修改 FIFO 替换策略或手动指定某次命中。

设关键再读 $e$ 的发射时刻为 $a_e$，其最近一次缓存填充为 $f_x$，中间将 $x$ 淘汰的事件为 $v_x$。一个看似自然的操作是让消费 $e$ 前移到淘汰之前；但这可能压住同核的矩阵/向量计算，或把新的 DDR 请求放进更拥堵的窗口。另一操作是把 $v_x$ 的相关子图后移，让 $x$ 暂时保留；这又可能推迟最终输出。故缓存账本只负责\emph{发现候选区域}，并用是否满足原 DAG、跨核 COPY 释放、Pipe 顺序及 L1/UB 容量来过滤提案。最终仍由完整官方 L2 评价计算候选的 $T^{L2}$ 和 $D$。

每个配置至多评估四个\emph{新增且不同}的候选，问题二种子与最终独立重放单独记账。若新候选均未改善实测字典序目标，则返回 $p_B$ 在 L2 下的评价结果；若有改善，则返回已评价方案中的最佳者。因而相对固定 L2 种子有逐配置不退步保障，而相对\emph{无 L2} 的 Q2 时间没有普遍保障：同一方案一旦加入共享 FIFO 与双带宽竞争，事件次序可能变，五百组中确有小幅负例。

搜索选择只应用了事件引导的迁核与插入、普通关键度后备和有限官方比较。开发阶段增加候选预算以及更复杂的关键链“消费前移+淘汰源后移”组合，未达到全量采用条件，不能把其开发样本收益写成当前五百组成绩。特别是在九个开发配置中，事件引导与普通邻域同预算对照为二胜六平一负；把候选上限从四增至十二也未让这九组最终时间继续下降。它支持“本轮采用较小预算”的工程选择，但不是搜索已收敛或更大预算永不收益的证明。

\subsection{命中敏感性与排序风险}
把张量读入从 DDR 改由 L2 服务，除了使独占工作量从 $\lceil b/60\rceil$ 改为 $\lceil b/250\rceil$，还会改变其完成时刻、后继算子发射、并发传输集合以及后续 FIFO 填充与淘汰次序。可用一个条件性扰动链表达：
\begin{equation}\label{eq:q3-causal-chain}
(C,b_L,p)\ \longrightarrow\
(Q_j,h_e,r(e))_{e}\ \longrightarrow\
(s_e,f_e)_{e}\ \longrightarrow\
\mathcal H(p)\ \longrightarrow\ T(p).
\end{equation}
箭头表示事件规则的依赖，不等于每一环都是可独立操纵的变量。固定 $p$ 增加 $C$ 或 $b_L$ 时，第一步可能使更多读取命中或更早完成，却也可能改变后续淘汰和竞争，最终 $T$ 不保证逐参数单调。这与式\eqref{eq:q3-relaxed-envelope}的乐观界单调并不矛盾：界允许任选 $z$ 平衡两个池，真实 FIFO 不允许自由选择命中分配。

由此可见，在问题三里用“命中字节最多”“额外搬运最少”或“传输独占时间最低”直接选方案都可能误判。对同一候选池若要改善排序器，应该固定 Q2 种子、缓存初态、候选集合与评价预算，比较代理排名和官方 $T^{L2}$ 排名；记录前两名官方覆盖率及选择遗憾，而不是仅展示代理预测下降的候选。只有在这些诊断证明筛选更可靠后，扩大缓存邻域才有根据。本文仍以已核验的有限候选算法交付，不把这些后续想法冒称为已实现收益。
